\section{Cryptographic Auditing \& Tamper-Evident Hash Chaining}
\label{sec:audit_cryptography}

\subsection{Enterprise Compliance \& Non-Repudiation}
In regulated enterprise environments (governed by Sarbanes-Oxley [SOX], SOC 2 Type II, and ISO 27001), granting autonomous AI agents the authority to mutate financial ledgers introduces severe compliance and liability requirements. Standard database transaction logs are insufficient: an insider or malicious actor with database write permissions could alter historical log entries after an incident to obscure unauthorized actions.

CentraOperator provides **mathematically verifiable, tamper-evident audit chaining** using SHA-256 cryptographic linkage:

\begin{figure}[htbp]
  \centering
  \begin{tikzpicture}[
  scale=0.88, transform shape,
  node distance=1.2cm and 1.6cm,
  block/.style={rectangle, rounded corners=2mm, draw=centrablue, fill=centrabg, line width=1.1pt, minimum width=3.3cm, minimum height=3.6cm, align=left, font=\scriptsize\sffamily, drop shadow},
  tampered/.style={rectangle, rounded corners=2mm, draw=centrared, fill=red!5, line width=1.3pt, minimum width=3.3cm, minimum height=3.6cm, align=left, font=\scriptsize\sffamily, drop shadow},
  hashflow/.style={->, >=Stealth, line width=1.4pt, draw=centrateal},
  badflow/.style={->, >=Stealth, line width=1.4pt, draw=centrared, dashed}
]

% Block 0 (Genesis)
\node[block] (b0) {
  \textbf{\large Block \#0 (Genesis)}\\[3pt]
  \textbf{Seq:} 0\\[2pt]
  \textbf{Timestamp:} $T_0$\\[2pt]
  \textbf{Event:} \texttt{task\_created}\\[2pt]
  \textbf{Prev Hash:}\\
  \texttt{0000000000000...}\\[3pt]
  \textbf{Payload:}\\
  \texttt{\{"goal": "Acme..."\}}\\[3pt]
  \textbf{Entry Hash ($H_0$):}\\
  \texttt{\textbf{a7f1...92b4}}
};

% Block 1
\node[block, right=1.2cm of b0] (b1) {
  \textbf{\large Block \#1 (Policy)}\\[3pt]
  \textbf{Seq:} 1\\[2pt]
  \textbf{Timestamp:} $T_1$\\[2pt]
  \textbf{Event:} \texttt{policy\_cleared}\\[2pt]
  \textbf{Prev Hash:}\\
  \texttt{\textbf{a7f1...92b4}}\\[3pt]
  \textbf{Payload:}\\
  \texttt{\{"amount": 1200.0\}}\\[3pt]
  \textbf{Entry Hash ($H_1$):}\\
  \texttt{\textbf{3e99...88cd}}
};

% Block 2
\node[block, right=1.2cm of b1] (b2) {
  \textbf{\large Block \#2 (Action)}\\[3pt]
  \textbf{Seq:} 2\\[2pt]
  \textbf{Timestamp:} $T_2$\\[2pt]
  \textbf{Event:} \texttt{web\_submit}\\[2pt]
  \textbf{Prev Hash:}\\
  \texttt{\textbf{3e99...88cd}}\\[3pt]
  \textbf{Payload:}\\
  \texttt{\{"status": "POSTED"\}}\\[3pt]
  \textbf{Entry Hash ($H_2$):}\\
  \texttt{\textbf{c4b1...007e}}
};

% Block 3
\node[block, right=1.2cm of b2] (b3) {
  \textbf{\large Block \#3 (Verify)}\\[3pt]
  \textbf{Seq:} 3\\[2pt]
  \textbf{Timestamp:} $T_3$\\[2pt]
  \textbf{Event:} \texttt{api\_verified}\\[2pt]
  \textbf{Prev Hash:}\\
  \texttt{\textbf{c4b1...007e}}\\[3pt]
  \textbf{Payload:}\\
  \texttt{\{"verified": true\}}\\[3pt]
  \textbf{Entry Hash ($H_3$):}\\
  \texttt{\textbf{f21a...4459}}
};

% Valid Hash Links
\draw[hashflow] (b0.east) -- (b1.west) node[midway, above, font=\tiny\bfseries\color{centrateal}] {$H_0$ Linked};
\draw[hashflow] (b1.east) -- (b2.west) node[midway, above, font=\tiny\bfseries\color{centrateal}] {$H_1$ Linked};
\draw[hashflow] (b2.east) -- (b3.west) node[midway, above, font=\tiny\bfseries\color{centrateal}] {$H_2$ Linked};

% Formula Annotation Below
\node[below=0.6cm of b1, xshift=2.0cm, align=center, font=\small\sffamily] (formula) {
  \color{centrablue}\textbf{Cryptographic Invariant:}\\
  $H_i = \text{SHA-256}\Big( H_{i-1} \parallel \text{canonical}\big(\text{entry\_body}_i\big) \Big)$\\
  \footnotesize Modifying any byte in block $k$ causes $H'_k \neq H_k$, deterministically invalidating all $i > k$.
};

\end{tikzpicture}

  \caption{Cryptographic Linkage in CentraOperator's Tamper-Evident Audit Chain. Every discrete state transition, action payload, gate decision, and verification check forms a sealed block cryptographically anchored to its predecessor.}
  \label{fig:audit_chain}
\end{figure}

\subsection{Formal Cryptographic Construction}
Let an audit chain $\mathcal{C}$ be an ordered sequence of entries $\langle e_0, e_1, \dots, e_n \rangle$. Each entry $e_i$ contains a payload body $B_i$, a sequence number $i$, a timestamp $t_i$, a predecessor hash $h_{i-1}$, and an entry hash $h_i$.

\begin{auditbox}[Formal Hash Chain Specification]
\begin{align}
  h_{\text{genesis}} &= 0^{64} = \underbrace{\texttt{0000000000000000000000000000000000000000000000000000000000000000}}_{64\text{ hexadecimal zeros}} \\
  h_0 &= \text{SHA-256}\Big( h_{\text{genesis}} \mathbin{\Vert} \text{canonical}(B_0) \Big) \\
  h_i &= \text{SHA-256}\Big( h_{i-1} \mathbin{\Vert} \text{canonical}(B_i) \Big), \quad \forall i \ge 1
\end{align}
where $\mathbin{\Vert}$ denotes byte concatenation, and $\text{canonical}(B)$ represents deterministic canonical JSON serialization.
\end{auditbox}

\subsubsection{Deterministic Canonical JSON Serialization}
Standard JSON serialization libraries produce non-deterministic strings due to arbitrary dictionary key orderings and varying whitespace formatting. CentraOperator enforces strict canonical serialization:
\begin{lstlisting}[language=Python, caption={Deterministic Canonical Serialization (\texttt{services/runtime/audit.py})}]
def canonical_json(obj: Any) -> str:
    """Produces deterministic canonical JSON serialization for cryptographic hashing."""
    return json.dumps(obj, sort_keys=True, separators=(',', ':'))
\end{lstlisting}
By enforcing lexicographical key sorting and eliminating extraneous whitespace separators, the hashed byte representation is strictly invariant across runtime environments, architectures, and language runtimes.

\subsection{Tamper Detection Mathematical Invariant}
Suppose an adversary modifies any field within an entry $e_k$ where $0 \le k < n$, replacing $B_k$ with $B'_k \ne B_k$.
\begin{enumerate}
  \item Due to the collision resistance of SHA-256:
  \begin{equation}
    h'_k = \text{SHA-256}\big(h_{k-1} \mathbin{\Vert} \text{canonical}(B'_k)\big) \ne h_k
  \end{equation}
  \item In the subsequent entry $e_{k+1}$, the stored predecessor hash remains $h_k$.
  \item When the verification routine evaluates block $k+1$:
  \begin{equation}
    e_{k+1}.\text{prev\_hash} = h_k \ne h'_k
  \end{equation}
  \item The chain verification halts immediately, pinpointing the exact block index of the integrity breach:
\end{enumerate}

\begin{lstlisting}[language=Python, caption={Chain Integrity Verification Routine (\texttt{services/runtime/audit.py})}]
def verify_chain_integrity(self) -> Dict[str, Any]:
    expected_prev = self.GENESIS_HASH
    for idx, entry in enumerate(self.chain):
        if entry["prev_hash"] != expected_prev:
            return {"valid": False, "tampered_at_index": idx, "reason": "Previous hash mismatch"}
        
        computed = hashlib.sha256(f"{expected_prev}{canonical_json(entry['body'])}".encode()).hexdigest()
        if computed != entry["hash"]:
            return {"valid": False, "tampered_at_index": idx, "reason": "Hash recalculation mismatch"}
        expected_prev = entry["hash"]
    return {"valid": True, "entries_count": len(self.chain), "final_hash": expected_prev}
\end{lstlisting}

This algorithm was rigorously validated in our automated test suite (\texttt{tests/test\_audit\_chain.py}), proving that tampering with any historical parameter (such as modifying an invoice amount from $\$1,200.00$ to $\$99,999.00$) causes an immediate integrity failure.

\subsection{Defense-in-Depth Transaction Idempotency}
In distributed finance architectures, network timeouts frequently create ambiguity: did a POST request time out before reaching the server, or did it succeed on the server while the response was lost in transit? Retrying naively risks duplicate ledger postings.

CentraOperator implements **three-layer defense-in-depth idempotency**:

\begin{enumerate}[leftmargin=*]
  \item \textbf{Deterministic Cryptographic Key Derivation:} Before issuing an external mutation, the agent computes an idempotency fingerprint:
  \begin{equation}
    K_{\text{idem}} = \text{SHA-256}\big(\text{operation} \mathbin{\Vert} \text{vendor} \mathbin{\Vert} \text{invoice\_number}\big)
  \end{equation}
  \item \textbf{Distributed Redis Lock with TTL:} The agent attempts an atomic \texttt{SET K\_idem "LOCKED" NX EX 60} in Redis. If the key already exists, the mutation is immediately aborted.
  \item \textbf{CentraERP Database Unique Composite Index:} The underlying datastore enforces a strict unique constraint on \texttt{(vendor, invoice\_number)}. Any concurrent write is rejected at the database engine level with an HTTP 409 Conflict.
  \item \textbf{Read-Before-Retry Protocol:} If an ambiguous network error occurs, the agent queries the read-only ERP API first to inspect whether the transaction landed before attempting a re-submission.
\end{enumerate}
